% Reviewed translation: PDF pages 251--255 / printed pages 237--241.
% The original scan is authoritative; mathematical tokens were checked against it.
\MathBookSourcePage{251}
此外, 假设 $\Omega$ 是凸的, 则问题
\eqref{eq:incompressible-sfvp-practical-problem} 的解
$\mathbf u_h=(\operatorname{curl}\psi_h,\omega_h)$ 满足
\begin{equation}\label{eq:incompressible-quadrilateral-velocity-superconvergence}
\|\mathbf u-\mathbf u_h\|_{\widetilde X}
\leq Ch^l\|\psi\|_{l+2,\Omega},
\end{equation}
其中使用另一个常数 $C>0$, 并假设 Stokes 问题
\eqref{eq:incompressible-sfvp-stokes} 的流函数
$\psi\in H^{l+2}(\Omega)$.
当 $\Omega$ 不凸时, 若涡量
$\omega\in H^{l+1}(\Omega)$, 仍可达到相同精度阶.

最后, 若问题 \eqref{eq:incompressible-sfvp-stokes}
的解 $(\mathbf u=(\operatorname{curl}\psi,\omega),p)$ 满足
$\psi\in H^{l+2}(\Omega)$, $p\in H^l(\Omega)$,
且 $\Omega$ 为凸域, 则问题
\eqref{eq:incompressible-sfvp-discrete-pressure} 的压力解满足
\begin{equation}\label{eq:incompressible-quadrilateral-pressure-superconvergence}
\|p-p_h\|_{0,\Omega}
\leq Ch^l\left(|p|_{l,\Omega}+\|\psi\|_{l+2,\Omega}\right).
\end{equation}

\section{二维 "流函数--速度梯度张量" 方法}
\label{sec:incompressible-stream-gradient-tensor}

本节讨论的方法以流函数 $\psi$ 和速度梯度
\[
\lambda_{ij}=\frac{\partial^2\psi}{\partial x_i\partial x_j},
\qquad1\leq i,j\leq2,
\]
作为未知量. 换言之, 它求解流函数及其全部二阶导数,
而不是像前一种方法那样只求 Laplacian.
读者将看到, 这种方法得到一种经济且最优的格式,
称为 Hellan--Herrmann--Johnson 格式.
Brezzi \& Raviart~\cite{ref:brezzi-raviart-15}
详尽分析了这一格式及相关格式.
这里给出的分析是 Babuška et al.~\cite{ref:babuska-osborn-pitkaranta-5}
提出的一种略有不同的分析, 它优雅地应用了依赖于网格的范数.

\subsection{Hellan--Herrmann--Johnson 表述}
\label{subsec:incompressible-hhj-formulation}

令 $(\mathbf u,p)$ 是问题
\eqref{eq:incompressible-sfvp-stokes} 的解.
不引入涡量 $\omega$, 而先把非对称速度梯度张量
$\boldsymbol\lambda=(\lambda_{ij})$
作为新的未知量:
\begin{equation}\label{eq:incompressible-hhj-velocity-gradient}
\lambda_{ij}=\frac{\partial u_i}{\partial x_j},
\qquad1\leq i,j\leq2.
\end{equation}
于是问题 \eqref{eq:incompressible-sfvp-stokes} 的方程可写为
\begin{equation}\label{eq:incompressible-hhj-gradient-system}
\left\{
\begin{aligned}
-\nu\sum_{j=1}^2\frac{\partial\lambda_{ij}}{\partial x_j}
+\frac{\partial p}{\partial x_i}&=f_i,
&&i=1,2,\\
\sum_{i=1}^2\lambda_{ii}&=0,
\end{aligned}
\right.
\quad\text{in }\Omega.
\end{equation}
将通过适当的分部积分从中导出新表述.

\MathBookSourcePage{252}
更精确地, 令 $\kappa$ 是 $\mathbb R^2$ 中有界开集,
其边界 $\partial\kappa$ Lipschitz 连续,
外单位法向量 $\mathbf n=(n_1,n_2)$,
单位切向量 $\mathbf t=(-n_2,n_1)$.
对 $\mathbf v\in H^1(\kappa)^2$
和张量 $\boldsymbol\tau=(\tau_{ij})\in H^1(\kappa)^4$,
Green 公式给出
\begin{equation}\label{eq:incompressible-hhj-green-component}
\int_\kappa\frac{\partial v_i}{\partial x_j}\tau_{ij}\,dx
=-\int_\kappa v_i\frac{\partial\tau_{ij}}{\partial x_j}\,dx
+\int_{\partial\kappa}v_i n_j\tau_{ij}\,ds.
\end{equation}
利用 $v_i=(\mathbf v\cdot\mathbf n)n_i
+(\mathbf v\cdot\mathbf t)t_i$, $i=1,2$,
边界积分也可写成
\[
\int_{\partial\kappa}v_i n_j\tau_{ij}\,ds
=\int_{\partial\kappa}(\mathbf v\cdot\mathbf n)
\tau_{ij}n_jn_i\,ds
+\int_{\partial\kappa}(\mathbf v\cdot\mathbf t)
\tau_{ij}n_jt_i\,ds.
\]
定义
\begin{equation}\label{eq:incompressible-hhj-normal-moments}
M_n(\boldsymbol\tau)
=\sum_{i,j=1}^2\tau_{ij}n_in_j,
\qquad
M_{nt}(\boldsymbol\tau)
=\sum_{i,j=1}^2\tau_{ij}n_jt_i,
\end{equation}
则对 $i,j$ 求和后,
\eqref{eq:incompressible-hhj-green-component} 化为
\[
\int_\kappa\frac{\partial v_i}{\partial x_j}\tau_{ij}\,dx
=-\int_\kappa v_i\frac{\partial\tau_{ij}}{\partial x_j}\,dx
+\int_{\partial\kappa}(\mathbf v\cdot\mathbf n)
M_n(\boldsymbol\tau)\,ds
+\int_{\partial\kappa}(\mathbf v\cdot\mathbf t)
M_{nt}(\boldsymbol\tau)\,ds.
\]
采用重复指标表示求和的惯例, 更紧凑地写成
\begin{equation}\label{eq:incompressible-hhj-green-tensor}
\begin{split}
\int_\kappa\frac{\partial v_i}{\partial x_j}\tau_{ij}\,dx
={}&-\int_\kappa v_i\frac{\partial\tau_{ij}}{\partial x_j}\,dx
+\int_{\partial\kappa}(\mathbf v\cdot\mathbf n)
M_n(\boldsymbol\tau)\,ds\\
&+\int_{\partial\kappa}(\mathbf v\cdot\mathbf t)
M_{nt}(\boldsymbol\tau)\,ds,
\end{split}
\end{equation}
对所有 $\mathbf v\in H^1(\kappa)^2$
和 $\boldsymbol\tau\in H^1(\kappa)^4$ 成立.
恒等式 \eqref{eq:incompressible-hhj-green-tensor}
与 \eqref{eq:incompressible-hhj-gradient-system}
共同构成包括 Hellan--Herrmann--Johnson 格式在内的若干格式的基础.

这里同样有多种方法表述当前方法.
我们采用与 Subsection~\ref{subsec:incompressible-sfvp-mesh-dependent}
密切相关的依赖于网格的表述.
设 $\Omega$ 是 $\mathbb R^2$ 中边界 $\Gamma$ 为多边形的有界区域,
$\mathcal T_h$ 是由三角形和/或凸四边形组成的剖分,
其单元直径由 $h$ 控制.
与 Subsection~\ref{subsec:incompressible-sfvp-mesh-dependent} 一样,
假设速度和右端项满足
$\mathbf u\in H^2(\Omega)^2$, $\mathbf f\in L^2(\Omega)^2$.

\MathBookSourcePage{253}
结合 \eqref{eq:incompressible-hhj-velocity-gradient}
和 \eqref{eq:incompressible-hhj-green-tensor}
考察方程 \eqref{eq:incompressible-hhj-gradient-system}.
若张量 $\boldsymbol\tau$ 整体属于 $H^1(\Omega)^4$,
而向量场 $\mathbf v$ 属于
\[
H_0(\operatorname{div};\Omega)
=\{\mathbf v\in L^2(\Omega)^2;
\ \operatorname{div}\mathbf v\in L^2(\Omega),
\ \mathbf v\cdot\mathbf n|_\Gamma=0\},
\]
则曲面积分之和
$\sum_{\kappa\in\mathcal T_h}
\int_{\partial\kappa}\mathbf v\cdot\mathbf n
M_n(\boldsymbol\tau)\,ds$ 为零.
因此 $\mathbf u\in H^2(\Omega)^2$ 的假设允许把
\eqref{eq:incompressible-hhj-green-tensor}
用于 \eqref{eq:incompressible-hhj-gradient-system} 的第一组方程, 得
\begin{equation}\label{eq:incompressible-hhj-velocity-pressure-weak}
\begin{split}
\nu\sum_{\kappa\in\mathcal T_h}
\biggl\{&\int_\kappa\lambda_{ij}
\frac{\partial v_i}{\partial x_j}\,dx
-\int_{\partial\kappa}M_{nt}(\boldsymbol\lambda)
\mathbf v\cdot\mathbf t\,ds\biggr\}
-\int_\Omega p\operatorname{div}\mathbf v\,dx\\
&=\int_\Omega\mathbf f\cdot\mathbf v\,dx,
\end{split}
\end{equation}
对所有 $\mathbf v\in H_0(\operatorname{div};\Omega)$,
$\mathbf v|_\kappa\in H^1(\kappa)^2$
于所有 $\kappa\in\mathcal T_h$ 成立.
同样, 为与 \eqref{eq:incompressible-hhj-velocity-pressure-weak} 匹配,
速度与梯度的关系可写成
\begin{equation}\label{eq:incompressible-hhj-gradient-relation-weak}
\int_\Omega\lambda_{ij}\tau_{ij}\,dx
=\sum_{\kappa\in\mathcal T_h}
\left\{
\int_\kappa\tau_{ij}\frac{\partial u_i}{\partial x_j}\,dx
-\int_{\partial\kappa}M_{nt}(\boldsymbol\tau)
\mathbf u\cdot\mathbf t\,ds
\right\}
\end{equation}
对所有 $\boldsymbol\tau\in H^1(\Omega)^4$ 成立,
因为对 $\mathbf u\in H_0^1(\Omega)^2$
和 $\boldsymbol\tau\in H^1(\Omega)^4$,
边界项之和为零.

观察 \eqref{eq:incompressible-hhj-velocity-pressure-weak}
和 \eqref{eq:incompressible-hhj-gradient-relation-weak} 可见,
测试函数 $\mathbf v$ 的正则性几乎不能降低,
但测试张量 $\boldsymbol\tau$ 不必整体属于 $H^1(\Omega)^4$.
事实上, 若 $\boldsymbol\tau|_\kappa\in H^1(\kappa)^4$
于每个 $\kappa$ 上, 则
\eqref{eq:incompressible-hhj-gradient-relation-weak} 已有意义.
此外, 方便假设 $M_n(\boldsymbol\tau)$
跨过所有单元间边界连续.

综上, 若 Stokes 问题 \eqref{eq:incompressible-sfvp-stokes}
的解 $(\mathbf u,p)$ 满足
$\mathbf u\in(H^2(\Omega)\cap H_0^1(\Omega))^2$,
$p\in H^1(\Omega)\cap L_0^2(\Omega)$,
则 $(\mathbf u,\boldsymbol\lambda,p)$,
$\lambda_{ij}=\partial u_i/\partial x_j$, 也是下列问题的解:
\begin{equation}\label{eq:incompressible-hhj-three-field-problem}
\left\{
\begin{aligned}
&\mathbf u\in H_0(\operatorname{div};\Omega),
\quad\mathbf u|_\kappa\in H^1(\kappa)^2,\\
&\boldsymbol\lambda|_\kappa\in H^1(\kappa)^4,
\quad M_n(\boldsymbol\lambda)\text{ 在 }\Gamma_h\text{ 上连续},
\quad p\in H^1(\Omega)\cap L_0^2(\Omega),\\
&\nu\sum_{\kappa\in\mathcal T_h}
\left\{\int_\kappa\lambda_{ij}\frac{\partial v_i}{\partial x_j}\,dx
-\int_{\partial\kappa}M_{nt}(\boldsymbol\lambda)
\mathbf v\cdot\mathbf t\,ds\right\}
-\int_\Omega p\operatorname{div}\mathbf v\,dx
=\int_\Omega\mathbf f\cdot\mathbf v\,dx,\\
&\int_\Omega\lambda_{ij}\tau_{ij}\,dx
-\sum_{\kappa\in\mathcal T_h}
\left\{\int_\kappa\tau_{ij}\frac{\partial u_i}{\partial x_j}\,dx
-\int_{\partial\kappa}M_{nt}(\boldsymbol\tau)
\mathbf u\cdot\mathbf t\,ds\right\}=0.
\end{aligned}
\right.
\end{equation}
其中第三行对所有分片 $H^1$ 的
$\mathbf v\in H_0(\operatorname{div};\Omega)$ 成立,
第四行对所有分片 $H^1$ 且 $M_n(\boldsymbol\tau)$
在 $\Gamma_h$ 上连续的张量成立.

反过来, 常规论证表明问题
\eqref{eq:incompressible-hhj-three-field-problem}
至多有一个解. 事实上, 第二组方程蕴含
$\lambda_{ij}=\partial u_i/\partial x_j$
于每个单元中, 进而 $\operatorname{div}\mathbf u=0$ 于 $\Omega$ 中.

\MathBookSourcePage{254}
若再令 $\mathbf f=0$, 容易得到所有 $\lambda_{ij}=0$.
于是 $\mathbf u$ 在每个单元中为常向量, 且
\[
\sum_{\kappa\in\mathcal T_h}
\int_{\partial\kappa}M_{nt}(\boldsymbol\tau)
\mathbf u\cdot\mathbf t\,ds=0
\]
对所有分片 $H^1$ 且 $M_n(\boldsymbol\tau)$
跨 $\Gamma_h$ 连续的张量成立.
结合 $\mathbf u\in H_0(\operatorname{div};\Omega)$,
立即得到 $\mathbf u=0$ 于 $\Omega$ 中.
所以只要原 Stokes 解具有足够正则性,
问题 \eqref{eq:incompressible-hhj-three-field-problem}
就是其等价表述.

为消去压力, 也可用流函数表述该问题.
回顾 $\mathbf v\in H_0(\operatorname{div};\Omega)$
满足 $\operatorname{div}\mathbf v=0$ 当且仅当
\[
\mathbf v=\operatorname{curl}\phi,
\qquad
\phi\in\Phi
=\{\phi\in H^1(\Omega);
\ \phi|_{\Gamma_0}=0,
\ \phi|_{\Gamma_i}\text{ 为常数},
\ 1\leq i\leq p\}.
\]
若此外 $\mathbf v|_\kappa\in H^1(\kappa)^2$,
则 $\phi|_\kappa\in H^2(\kappa)$, 反之亦然.
可以直接用流函数改写
\eqref{eq:incompressible-hhj-three-field-problem},
但若不使用 $\boldsymbol\lambda$, 而引入对称张量
\begin{equation}\label{eq:incompressible-hhj-symmetric-tensor}
\sigma_{11}=-\lambda_{12},
\qquad
\sigma_{12}=\sigma_{21}=\lambda_{11},
\qquad
\sigma_{22}=\lambda_{21},
\end{equation}
记号会更简单. $\boldsymbol\sigma$ 与
$\boldsymbol\lambda$ 的对应关系为
\begin{equation}\label{eq:incompressible-hhj-tensor-correspondence}
\lambda_{11}=-\sigma_{12},
\quad\lambda_{12}=-\sigma_{22},
\quad\lambda_{21}=\sigma_{11},
\quad\lambda_{22}=\sigma_{12},
\quad
M_{nt}(\boldsymbol\lambda)=-M_n(\boldsymbol\sigma),
\quad
M_n(\boldsymbol\lambda)=M_{nt}(\boldsymbol\sigma).
\end{equation}
因此定义张量空间
\begin{equation}\label{eq:incompressible-hhj-tensor-space}
\mathcal X
=\{\boldsymbol\tau\in L^2(\Omega)^4;
\ \boldsymbol\tau|_\kappa\in H^1(\kappa)^4,
\ \tau_{ij}=\tau_{ji},
\ M_n(\boldsymbol\tau)\text{ 在 }\Gamma_h\text{ 的每条边上连续}\},
\end{equation}
以及 Subsection~\ref{subsec:incompressible-sfvp-mesh-dependent}
已经引入的流函数空间
\begin{equation}\label{eq:incompressible-hhj-stream-space}
\widetilde\Psi
=\{\phi\in\Phi;
\ \phi|_\kappa\in H^2(\kappa)
\ \forall\kappa\in\mathcal T_h\}.
\end{equation}
借助上述对应关系,
\eqref{eq:incompressible-hhj-three-field-problem} 中曲面积分之和可写成
\[
\sum_{\kappa\in\mathcal T_h}
\int_{\partial\kappa}M_{nt}(\boldsymbol\lambda)
\mathbf u\cdot\mathbf t\,ds
=\int_{\Gamma_h}M_n(\boldsymbol\sigma)
S(\partial\psi/\partial n)\,ds,
\]
其中 $S(v)$ 表示 $v$ 跨 $\Gamma_h$ 各边的跳跃.
因此得到称为 Hellan--Herrmann--Johnson 表述的等价问题:

\MathBookSourcePage{255}
求 $(\boldsymbol\sigma,\psi)\in\mathcal X\times\widetilde\Psi$, 使
\begin{equation}\label{eq:incompressible-hhj-formulation}
\left\{
\begin{aligned}
\nu\sum_{\kappa\in\mathcal T_h}
\int_\kappa\frac{\partial^2\phi}{\partial x_i\partial x_j}
\sigma_{ij}\,dx
-\nu\int_{\Gamma_h}M_n(\boldsymbol\sigma)
S(\partial\phi/\partial n)\,ds
&=\int_\Omega\mathbf f\cdot\operatorname{curl}\phi\,dx
&&\forall\phi\in\widetilde\Psi,\\
\int_\Omega\sigma_{ij}\tau_{ij}\,dx
-\sum_{\kappa\in\mathcal T_h}
\int_\kappa\frac{\partial^2\psi}{\partial x_i\partial x_j}
\tau_{ij}\,dx
+\int_{\Gamma_h}M_n(\boldsymbol\tau)
S(\partial\psi/\partial n)\,ds
&=0
&&\forall\boldsymbol\tau\in\mathcal X.
\end{aligned}
\right.
\end{equation}
再次可用两个双线性形式表述它:
\begin{equation}\label{eq:incompressible-hhj-form-a}
a_h(\boldsymbol\sigma,\boldsymbol\tau)
=\int_\Omega\sigma_{ij}\tau_{ij}\,dx
\qquad\forall\boldsymbol\sigma,\boldsymbol\tau\in L^2(\Omega)^4,
\end{equation}
以及
\begin{equation}\label{eq:incompressible-hhj-form-b}
b_h(\boldsymbol\tau,\phi)
=-\sum_{\kappa\in\mathcal T_h}
\int_\kappa\frac{\partial^2\phi}{\partial x_i\partial x_j}
\tau_{ij}\,dx
+\int_{\Gamma_h}M_n(\boldsymbol\tau)
S(\partial\phi/\partial n)\,ds.
\end{equation}
于是问题 \eqref{eq:incompressible-hhj-formulation} 写成
\begin{equation}\label{eq:incompressible-hhj-abstract-formulation}
\left\{
\begin{aligned}
b_h(\boldsymbol\sigma,\phi)
&=-\frac1\nu(\mathbf f,\operatorname{curl}\phi)
&&\forall\phi\in\widetilde\Psi,\\
a_h(\boldsymbol\sigma,\boldsymbol\tau)
+b_h(\boldsymbol\tau,\psi)&=0
&&\forall\boldsymbol\tau\in\mathcal X.
\end{aligned}
\right.
\end{equation}
注意它与 Subsection~\ref{subsec:incompressible-sfvp-mesh-dependent}
的问题 \eqref{eq:incompressible-sfvp-mesh-discrete-problem}
和 \eqref{eq:incompressible-sfvp-mesh-abstract-problem} 的类比.
还要注意, 当 $\boldsymbol\tau\in H^1(\Omega)^4$
且 $\phi\in\widetilde\Psi$ 时,
\begin{equation}\label{eq:incompressible-hhj-form-b-gradient}
b_h(\boldsymbol\tau,\phi)
=\int_\Omega
\frac{\partial\phi}{\partial x_i}
\frac{\partial\tau_{ij}}{\partial x_j}\,dx,
\end{equation}
这一性质类似于
\eqref{eq:incompressible-sfvp-mesh-form-b-curl}.

\begin{Theorem}\label{thm:incompressible-hhj-continuous-equivalence}
令 $\Omega$ 是有界平面多边形,
$\mathcal T_h$ 是 $\Omega$ 的一个剖分.
假设 Stokes 问题 \eqref{eq:incompressible-sfvp-stokes}
的解 $(\mathbf u=\operatorname{curl}\psi,p)$ 满足
\[
\mathbf u\in(H^2(\Omega)\cap H_0^1(\Omega))^2,
\qquad p\in H^1(\Omega)\cap L_0^2(\Omega).
\]
则:

\textup{1)} 三元组
$(\mathbf u,(\lambda_{ij}=\partial u_i/\partial x_j),p)$
是问题 \eqref{eq:incompressible-hhj-three-field-problem} 的唯一解;

\textup{2)} 二元组
$((\sigma_{ij}=\partial^2\psi/
\partial x_i\partial x_j),\psi)$
是问题 \eqref{eq:incompressible-hhj-formulation} 的唯一解.
\end{Theorem}

下一小节将假设
\begin{equation}\label{eq:incompressible-hhj-regularity-assumption}
\psi\in H^3(\Omega)\cap H_0^2(\Omega),
\end{equation}
以便可以使用问题
\eqref{eq:incompressible-hhj-three-field-problem}
或 \eqref{eq:incompressible-hhj-formulation} 中的任一个.
